\section*{Appendix D: A Lean Learning Experiment and Its Decision Gates}

The following 30-day sprint is designed for a founder with limited resources. The time blocks are illustrative operational guidance, not experimental benchmarks.

\textbf{Days 1 to 5: Select 1 capability and diagnose the baseline.} Talk with real practitioners about costly failures and collect privacy-safe cases. Define the relevant tasks and assistance conditions. Give prospective learners 1 unseen, appropriately scoped diagnostic. Analyze reasoning and error patterns. Reject or narrow a target that cannot be observed with available methods.

\textbf{Days 6 to 10: Build the minimal learning loop.} Develop 1 accurate explanation, 2 contrasts, 1 worked case, a guided task, an independent task, a feedback rubric, and a retry. Reserve later tasks that will not appear in the teaching materials. Plan how to handle a learner whose prerequisites are insufficient.

\textbf{Days 11 to 18: Teach a small pilot.} Observe individual attempts, not merely participation. Deliver promised feedback and require revisions. Note what human reviewers do that automated checks cannot reliably perform. Record elapsed time and support burden. Ensure participants understand what is experimental and what is guaranteed.

\textbf{Days 19 to 24: Challenge independent performance.} Use an unseen, representative task scored against predeclared criteria. Include an unfamiliar variant if transfer is part of the promise. Record tool use and outside assistance. Do not manufacture a success rate by excluding those who left the program.

\textbf{Days 25 to 30: Decide.} Compare entry and exit performance cautiously, review error types, obtain learner reports of value, and calculate true delivery effort. Choose among \emph{repeat with improvements}, \emph{narrow the skill or learner segment}, \emph{change the learning mechanism}, and \emph{stop}. A delayed check, if feasible, should be scheduled with consent; immediate results are not the same as long-term retention.

\subsection*{D1. Predeclared Decision Questions}

\begin{enumerate}
    \item Did learners demonstrate a meaningful improvement in the target task family?
    \item Did improvements survive a task not directly rehearsed?
    \item Were the relevant errors diagnosed and repaired rather than hidden?
    \item Can learners realistically sustain the required practice schedule?
    \item Did the feedback workload remain within the delivery capacity promised?
    \item Was the experience valuable enough for a realistic buyer to commit resources?
    \item What alternative explanation could account for the apparent gain?
    \item What single design change should the next cohort test?
\end{enumerate}

A strong 1st pilot may justify a repeat. It does not establish product-market fit, causal superiority, or permanent career transformation.

\section*{Appendix E: The Educational Founder's Evidence Register}

Keep the distinction between \textbf{what the institution did} and \textbf{what the learner became able to do} visible.

\begin{table*}[t]
\centering
\small
\caption{The Educational Founder's Evidence Register}
\begin{tabularx}{\textwidth}{p{0.18\textwidth}p{0.26\textwidth}p{0.26\textwidth}X}
\toprule
\textbf{Type of Evidence} & \textbf{Examples} & \textbf{What It Legitimately Supports} & \textbf{What It Cannot Establish Alone} \\
\midrule
Exposure & Hours, video completion, attendance & Opportunity to learn & Mastery \\
Reaction & Ratings, enjoyment, confidence & Experienced value and engagement & Accurate competence \\
Near performance & Immediate quizzes and rehearsed variants & Short-run understanding or procedure & Retention and transfer \\
Independent performance & Unseen task under specified assistance & Capability on sampled tasks & Generalization to all work \\
Delayed performance & Similar competence after time passes & Some durability & Full professional expertise \\
Transfer performance & Changed contexts without announced method & Bounded generalization & Universal problem-solving ability \\
Authentic work & Verified professional artifacts and decisions & Relevance to work domain & That instruction caused the result \\
Comparative evidence & Credibly designed comparison of approaches & Stronger attribution in tested setting & Universal applicability \\
Commercial evidence & Payments, referrals, renewals & Buyer demand and perceived value & Educational effectiveness \\
\bottomrule
\end{tabularx}
\end{table*}

\textbf{A disciplined report states the denominator.} ``9 of 12 enrolled learners demonstrated the standard on a defined final assessment'' conveys more than ``75\% success'' when missing participants and assessment conditions are hidden. If the 9 were the only ones assessed, the claim must instead say so. Where no comparable baseline exists, refrain from claiming an improvement caused by instruction.

The reporting objective is not to make the company sound cautious for its own sake. It is to make its educational claims useful and credible to learners, parents, employers, partners, and eventually independent academic reviewers.

\section*{Appendix F: 10 Distinctions That Should Govern Institutional Decisions}

\textbf{1. Learning versus performance.} Producing an answer with support does not necessarily imply knowledge that persists when support is removed.

\textbf{2. Comprehension versus independent construction.} Understanding a demonstration is not the same as creating a solution from an unfamiliar problem.

\textbf{3. Recognition versus retrieval.} Familiar content is not necessarily accessible when no prompt supplies it.

\textbf{4. Correct output versus correct reasoning.} A plausible answer can be obtained through a lucky guess, template, collaborator, or AI assistant.

\textbf{5. Difficulty versus useful difficulty.} A broken tool wastes effort; diagnosing a carefully chosen unfamiliar problem can build skill.

\textbf{6. Feedback versus repair.} Receiving comments does not imply the learner made a better subsequent attempt.

\textbf{7. Selectivity versus value added.} Excellent entrants can produce excellent graduates without large educational contribution.

\textbf{8. Retention versus transfer.} Remembering the original case is not equivalent to recognizing the principle in a different setting.

\textbf{9. Educational effectiveness versus commercial repeatability.} A powerful learning experience may be financially unviable; a popular course may teach little.

\textbf{10. Mission versus evidence.} Wanting to expand human potential is an aspiration. Demonstrating a bounded, repeatable capability gain is the institution's responsibility.

\section*{Appendix G: Questions for Faculty, Founders, and Learners}

\subsection*{For a Faculty Member Reviewing a Lesson}

What is the \textbf{single most consequential misconception} this lesson is designed to prevent? What prior knowledge makes the explanation intelligible? Which visible attempt tells us whether the lesson worked? What would a knowledgeable but incorrect student answer, and why? Is the feedback specific enough to change the next attempt? Where does the concept reappear after a delay? What evidence would convince you that the lesson needs to be redesigned?

\subsection*{For a Founder Deciding Whether to Fund a Course}

Which real buyer values this capability? What existing alternative are learners using? Which transformation can be demonstrated with the current team and constraints? Which expensive feedback work is unavoidable? Are there cheaper ways to supply foundational explanation without sacrificing quality? What will we refuse to promise? What would justify a 2nd cohort, and what would make us stop?

\subsection*{For a Learner Deciding Whether to Enroll}

Can I inspect the exact work I am expected to become capable of doing? Are the prerequisites stated clearly? Will I attempt difficult tasks independently? What kind of feedback will I actually receive, and when? Must I revise work after feedback? How will I know the difference between completing the course and meeting its standard? How much weekly effort will be required? What outcomes are explicitly \textbf{not} promised?

\subsection*{For the Institution Reviewing an Unsuccessful Graduate}

Did the learner begin with sufficient prerequisites? Were the expectations and workload feasible? Were errors detected early? Was help requested or offered in time? Did the feedback lead to a reattempt? Was the assessment representative and fair? Was the failure specific to a capability, or did the institution infer a broad deficiency from an unrelated barrier? What should change before the next delivery?

The questions are not a substitute for scientific research. They make the institute's causal theory inspectable by the people who will bear the consequences.
